\SourcePage{163}{166}
\section{Motion in a Coulomb field (parabolic coordinates)}

For motion in any centrally symmetric field, the Schrödinger equation can
always be separated in spherical coordinates. In a Coulomb field, separation
is also possible in the so-called parabolic coordinates. The solution of the
Coulomb-field problem in these coordinates is useful in studying a number of
problems in which a particular direction in space is singled out, for example
by an external electric field in addition to the Coulomb field (\secref{77}).
\SourcePage{164}{167}
The parabolic coordinates $\xi$, $\eta$, and $\varphi$ are defined by
\begin{equation}
 \begin{aligned}
  x&=\sqrt{\xi\eta}\cos\varphi,\qquad
  y=\sqrt{\xi\eta}\sin\varphi,\qquad
  z=\frac12(\xi-\eta),\\
  r&=\sqrt{x^2+y^2+z^2}=\frac12(\xi+\eta)
 \end{aligned}
 \tag{37.1}
\end{equation}
or, conversely,
\begin{equation}
  \xi=r+z,\qquad \eta=r-z,\qquad
  \varphi=\arctan\frac yx;
  \tag{37.2}
\end{equation}
$\xi$ and $\eta$ range from 0 to $\infty$, while $\varphi$ ranges from 0 to
$2\pi$. The surfaces $\xi=\operatorname{const}$ and
$\eta=\operatorname{const}$ are paraboloids of revolution whose axis is the
$z$ axis and whose focus is at the origin. This coordinate system is
orthogonal. The line element is
\begin{equation}
  dl^2=\frac{\xi+\eta}{4\xi}\,d\xi^2
       +\frac{\xi+\eta}{4\eta}\,d\eta^2
       +\xi\eta\,d\varphi^2,
  \tag{37.3}
\end{equation}
and the volume element is
\begin{equation}
  dV=\frac14(\xi+\eta)\,d\xi\,d\eta\,d\varphi.
  \tag{37.4}
\end{equation}
Equation (37.3) gives the following expression for the Laplace operator:
\begin{equation}
  \Delta=\frac4{\xi+\eta}\left[
    \frac\partial{\partial\xi}\left(\xi\frac\partial{\partial\xi}\right)
   +\frac\partial{\partial\eta}\left(\eta\frac\partial{\partial\eta}\right)
  \right]
  +\frac1{\xi\eta}\frac{\partial^2}{\partial\varphi^2}.
  \tag{37.5}
\end{equation}

For a particle in the attractive Coulomb field
\[
  U=-\frac1r=-\frac2{\xi+\eta},
\]
the Schrödinger equation becomes
\begin{equation}
  \frac4{\xi+\eta}\left[
    \frac\partial{\partial\xi}\left(\xi\frac{\partial\psi}{\partial\xi}\right)
   +\frac\partial{\partial\eta}\left(\eta\frac{\partial\psi}{\partial\eta}\right)
  \right]
  +\frac1{\xi\eta}\frac{\partial^2\psi}{\partial\varphi^2}
  +2\left(E+\frac2{\xi+\eta}\right)\psi=0.
  \tag{37.6}
\end{equation}
We seek the eigenfunctions $\psi$ in the form
\begin{equation}
  \psi=f_1(\xi)f_2(\eta)e^{im\varphi},
  \tag{37.7}
\end{equation}
where $m$ is the magnetic quantum number. Substitute this expression into
(37.6), after multiplying that equation by $(\xi+\eta)/4$, and separate the
variables $\xi$ and $\eta$. The resulting equations for $f_1$ and $f_2$ are
\begin{equation}
 \begin{aligned}
  \frac d{d\xi}\left(\xi\frac{df_1}{d\xi}\right)
  +\left[\frac E2\xi-\frac{m^2}{4\xi}+\beta_1\right]f_1&=0,\\
  \frac d{d\eta}\left(\eta\frac{df_2}{d\eta}\right)
  +\left[\frac E2\eta-\frac{m^2}{4\eta}+\beta_2\right]f_2&=0,
 \end{aligned}
 \tag{37.8}
\end{equation}
\SourcePage{165}{168}
where the ``separation parameters'' $\beta_1$ and $\beta_2$ are related by
\begin{equation}
  \beta_1+\beta_2=1.
  \tag{37.9}
\end{equation}

Consider the discrete energy spectrum ($E<0$). In place of $E$, $\xi$, and
$\eta$, introduce the quantities
\begin{equation}
  n=\frac1{\sqrt{-2E}},\qquad
  \rho_1=\xi\sqrt{-2E}=\frac\xi n,\qquad
  \rho_2=\frac\eta n,
  \tag{37.10}
\end{equation}
after which the equation for $f_1$ takes the form
\begin{equation}
  \frac{d^2f_1}{d\rho_1^2}+\frac1{\rho_1}\frac{df_1}{d\rho_1}
  +\left[-\frac14+\frac1{\rho_1}
    \left(\frac{|m|+1}{2}+n_1\right)-\frac{m^2}{4\rho_1^2}\right]f_1=0,
  \tag{37.11}
\end{equation}
with an identical equation for $f_2$; here we have also introduced the
notation
\begin{equation}
  n_1=-\frac{|m|+1}{2}+n\beta_1,\qquad
  n_2=-\frac{|m|+1}{2}+n\beta_2.
  \tag{37.12}
\end{equation}

Proceeding as for equation (36.4), we find that at large $\rho_1$, $f_1$
behaves as $e^{-\rho_1/2}$, and at small $\rho_1$ as
$\rho_1^{|m|/2}$. We therefore seek a solution of (37.11) in the form
\[
  f_1(\rho_1)=e^{-\rho_1/2}\rho_1^{|m|/2}w_1(\rho_1)
\]
(and analogously for $f_2$), obtaining for $w_1$ the equation
\[
  \rho_1w_1''+(|m|+1-\rho_1)w_1'+n_1w_1=0.
\]
This is again the confluent hypergeometric equation. The solution satisfying
the finiteness conditions is
\[
  w_1=F(-n_1,|m|+1,\rho_1),
\]
and $n_1$ must be a nonnegative integer.

Thus, in parabolic coordinates every stationary state of the discrete
spectrum is specified by three integers: the ``parabolic quantum numbers''
$n_1$ and $n_2$, and the magnetic quantum number $m$. From (37.9) and
(37.12), the number $n$ (the ``principal quantum number'') is
\begin{equation}
  n=n_1+n_2+|m|+1.
  \tag{37.13}
\end{equation}
The energy levels are, of course, the same as in (36.9).

For a given $n$, the number $|m|$ can take the $n$ values from 0 through
$n-1$. With $n$ and $|m|$ fixed, $n_1$ takes
\SourcePage{166}{169}
the $n-|m|$ values from 0 through $n-|m|-1$. For each prescribed $|m|$,
functions with $m=\pm|m|$ may also be chosen. Hence the total number of
distinct states for a given $n$ is
\[
  2\sum_{m=1}^{n-1}(n-m)+(n-0)=n^2,
\]
in agreement with the result obtained in \secref{36}.

The discrete-spectrum wave functions $\psi_{n_1n_2m}$ must obey the
normalization condition
\begin{equation}
  \int|\psi_{n_1n_2m}|^2\,dV
  =\frac14\int_0^\infty\!\int_0^\infty\!\int_0^{2\pi}
   |\psi_{n_1n_2m}|^2(\xi+\eta)\,d\varphi\,d\xi\,d\eta=1.
  \tag{37.14}
\end{equation}
The normalized functions are
\begin{equation}
  \psi_{n_1n_2m}=\frac{\sqrt2}{n^2}
  f_{n_1m}\left(\frac\xi n\right)
  f_{n_2m}\left(\frac\eta n\right)
  \frac{e^{im\varphi}}{\sqrt{2\pi}},
  \tag{37.15}
\end{equation}
\begin{equation}
  f_{pm}(\rho)=\frac1{|m|!}\sqrt{\frac{(p+|m|)!}{p!}}
  F(-p,|m|+1,\rho)e^{-\rho/2}\rho^{|m|/2}.
  \tag{37.16}
\end{equation}

Unlike the wave functions in spherical coordinates, those in parabolic
coordinates are not symmetric about the plane $z=0$. When $n_1>n_2$, the
particle is more likely to be found on the $z>0$ side than on the $z<0$
side; for $n_1<n_2$, the reverse holds.

The continuous spectrum ($E>0$) corresponds to a continuous spectrum of real
values of the parameters $\beta_1$, $\beta_2$ in equations (37.8) (still, of
course, related by (37.9)); we shall not write the corresponding wave
functions here. Regarded as equations for the ``eigenvalues'' $\beta_1$ and
$\beta_2$, equations (37.8) also possess, when $E>0$, a spectrum of complex
values. The associated wave functions will be given in \secref{135}, where they
will be used to solve the problem of scattering in a Coulomb field.

The existence of the stationary states $|n_1n_2m\rangle$ is tied to the
additional conservation law (36.29). In these states, besides the energy, the
quantities $l_z=m$ and $A_z$ have definite values. Evaluation of the diagonal
matrix elements of the operator $\hat A_z$ gives
\begin{equation}
  A_z=\frac{n_1-n_2}{n}.
  \tag{37.17}
\end{equation}
\SourcePage{167}{170}
Here $u_z=n_1-n_2$, while the projections of the ``angular momenta''
$\mathbf j_1$ and $\mathbf j_2$ are
\begin{equation}
  j_{1z}=\frac{m+n_1-n_2}{2}\equiv\mu_1,\qquad
  j_{2z}=\frac{m-n_1+n_2}{2}\equiv\mu_2.
  \tag{37.18}
\end{equation}

These properties of the states $|n_1n_2m\rangle$ (or, equivalently,
$|n\mu_1\mu_2\rangle$) make it easy to relate their wave functions to those
of the states $|nlm\rangle$. Since $\mathbf l=\mathbf j_1+\mathbf j_2$, a
change between the two descriptions reduces to constructing wave functions
for the addition of two angular momenta (considered below in \secref{106}). In terms
of the ``angular momenta'' $\mathbf j_1$ and $\mathbf j_2$, the states
$|nlm\rangle$ and $|n_1n_2m\rangle$ are written as
$|j_1j_2lm\rangle$ and $|j_1j_2\mu_1\mu_2\rangle$, where, by (36.35) and
(37.13),
\begin{equation}
  j_1=j_2=\frac{n-1}{2}=\frac{n_1+n_2+|m|}{2}.
  \tag{37.19}
\end{equation}
The general formulas (106.9)--(106.11) give
\begin{equation}
 \begin{aligned}
  \psi_{nlm}&=\sum_{\mu_1+\mu_2=m}
   \langle lm|\mu_1\mu_2\rangle\psi_{n\mu_1\mu_2},\\
  \psi_{n\mu_1\mu_2}&=\sum_{l=0}^{n-1}
   \langle l,\mu_1+\mu_2|\mu_1\mu_2\rangle\psi_{nlm}
 \end{aligned}
 \tag{37.20}
\end{equation}
(\emph{D. Park}, 1960).
